\chapter{Distributed Training Systems}
\label{ch:distributed-training-systems}

\abstract{This chapter covers the systems layer that makes large-scale training possible: data parallelism, tensor parallelism, pipeline parallelism, fully sharded training, ZeRO-style optimizer partitioning, activation checkpointing, NCCL communication, and throughput accounting.}

\section{Chapter Spine}
Training systems determine whether a model recipe is feasible. The chapter uses memory and bandwidth accounting as the bridge between algorithmic intent and hardware execution.

\section{Core Sections}
\subsection{Memory Accounting}
Parameters, gradients, optimizer states, activations, temporary buffers, and fragmentation.

\subsection{Parallelism Strategies}
DDP, FSDP, ZeRO, tensor parallelism, pipeline parallelism, expert parallelism, and hybrid layouts.

\subsection{Precision and Communication}
FP32, BF16, FP16, FP8, loss scaling, collectives, overlap, and network topology.

\subsection{Operational Reliability}
Checkpoint cadence, resumption, straggler detection, data-loader bottlenecks, and experiment tracking.

\section{Planned Exercises}
Compute per-GPU memory for a small decoder and explain which parallelism strategy fixes each bottleneck.
